\documentclass[11pt,letterpaper]{article}
\usepackage[top=0.7in,bottom=0.7in,left=0.75in,right=0.75in]{geometry}
\usepackage[T1]{fontenc}
\usepackage[utf8]{inputenc}
\usepackage{lmodern}
\usepackage{microtype}
\usepackage{enumitem}
\usepackage[hidelinks]{hyperref}
\usepackage{titlesec}
\setlength{\parindent}{0pt}
\setlength{\parskip}{0pt}
\pagestyle{empty}
\titleformat{\section}{\normalfont\bfseries}{}{0em}{\uppercase}[\vspace{-4pt}\hrule\vspace{4pt}]
\titlespacing{\section}{0pt}{8pt}{5pt}
\setlist[itemize]{leftmargin=1.2em,label=--,itemsep=0pt,topsep=2pt,parsep=0pt}
\begin{document}
\begin{center}
{\fontsize{16}{19}\selectfont\textbf{\MakeUppercase{Diego Castillo Pineda}}}\\[4pt]
{\small La Florida, Santiago, Chile \quad|\quad +56 9 94162547 \quad|\quad \href{mailto:diego.castillop11@gmail.com}{diego.castillop11@gmail.com}\\
\href{https://https://www.linkedin.com/in/diego-castillo-pineda/}{https://www.linkedin.com/in/diego-castillo-pineda/} \quad|\quad \href{https://https://github.com/diegocastillop11-cloud}{https://github.com/diegocastillop11-cloud}}
\end{center}

\section{Resumen Profesional}
Analista de Datos especializado en SQL Server, Python y visualización con Power BI, con experiencia comprobada en administración de bases de datos de alta concurrencia y optimización de procesos analíticos. Experto en extracción, transformación y análisis de datos mediante procedimientos almacenados, scripts automatizados y herramientas de Business Intelligence, logrando reducir tiempos de procesamiento en 40\% mediante automatización T-SQL en entornos con más de 50,000 transacciones por hora. Dominio de Azure SQL, Oracle y tecnologías cloud para integración de datos, con capacidad para traducir requerimientos de negocio en soluciones técnicas basadas en datos que impactan directamente en la toma de decisiones estratégicas.

\section{Experiencia Profesional}

\textbf{Punto Ticket} \hfill Santiago, Chile\\
\textit{Analista de Base de Datos Senior} \hfill Nov. 2019 – Feb. 2026
\begin{itemize}
  \item Diseñé y optimicé más de 120 procedimientos almacenados en SQL Server para análisis de datos transaccionales en eventos masivos, procesando más de 50,000 transacciones por hora y manteniendo tiempos de respuesta bajo 2 segundos durante períodos críticos de venta
  \item Automaticé procesos de extracción y transformación de datos mediante scripts T-SQL y Python, reduciendo errores manuales en 85\% y disminuyendo tiempos de generación de reportes analíticos de 4 horas a 30 minutos
  \item Desarrollé consultas dinámicas con funciones PIVOT y análisis multidimensional para reportería ejecutiva, entregando dashboards interactivos que facilitaron la toma de decisiones estratégicas en tiempo real
  \item Integré flujos de datos entre sistemas heterogéneos procesando archivos JSON, XML y APIs REST, asegurando consistencia y calidad de datos para análisis downstream
  \item Participé activamente en migración de bases de datos on-premise hacia Microsoft Azure SQL Database, optimizando arquitectura cloud para análisis escalable y reducción de costos operativos en 30\%
  \item Administré bases de datos SQL Server en alta disponibilidad, implementando monitoreo proactivo y optimización de índices que mejoraron performance de consultas analíticas en 45\%
\end{itemize}
\vspace{2pt}
\textbf{Imperial S.A.} \hfill Santiago, Chile\\
\textit{Analista Funcional} \hfill Jun. 2017 – Nov. 2019
\begin{itemize}
  \item Realicé análisis funcional y extracción de datos desde SQL Server y Oracle para resolver incidencias operativas, trabajando con más de 15 áreas de negocio para identificar patrones y oportunidades de mejora mediante análisis de datos
  \item Levanté requerimientos analíticos junto a usuarios clave, diseñando reportes personalizados en SQL Server Management Studio y Excel que automatizaron el seguimiento de KPIs críticos del negocio
  \item Parametricé y configuré módulos de sistemas ERP basándome en análisis de datos históricos, optimizando flujos de información y reduciendo tiempos de cierre contable en 25\%
  \item Desarrollé documentación técnica incluyendo diccionarios de datos, diagramas de flujo y manuales de usuario para facilitar la interpretación y uso de información analítica por equipos no técnicos
  \item Capacité a más de 50 usuarios finales en herramientas de consulta de datos y generación de reportes, promoviendo cultura data-driven dentro de la organización
\end{itemize}
\vspace{2pt}
\textbf{OB GROUP PARK} \hfill Santiago, Chile\\
\textit{Analista de Sistemas y TI} \hfill Mayo 2014 – Abril 2017
\begin{itemize}
  \item Lideré proyectos de integración tecnológica en retail, implementando soluciones de automatización que mejoraron la exactitud de datos contables en 35\% y redujeron tiempos de consolidación financiera
  \item Coordiné implementación de nuevas plataformas de gestión, asegurando migración exitosa de datos históricos y compatibilidad entre sistemas legacy y modernos sin pérdida de información crítica
  \item Colaboré con equipos multidisciplinarios para definir arquitectura de datos que soportara reportería operativa en tiempo real, mejorando visibilidad de indicadores clave de desempeño del negocio
\end{itemize}
\vspace{2pt}
\textbf{Hewlett-Packard Company} \hfill Santiago, Chile\\
\textit{Agente de Mesa de Ayuda Nivel 1} \hfill Ago. 2013 – Mayo 2014
\begin{itemize}
  \item Brindé soporte técnico a usuarios finales resolviendo más de 40 tickets semanales, documentando incidencias en sistema de gestión y generando bases de conocimiento que redujeron tiempos de resolución en 20\%
\end{itemize}

\section{Proyectos Destacados}

\textbf{Plataforma de Análisis Predictivo con IA y React} \hfill 2024
\begin{itemize}
  \item Desarrollé aplicación web full-stack combinando React (TypeScript) para frontend, Node.js + Express para backend, y PostgreSQL como base de datos, integrando APIs de Claude AI y Gemini para análisis predictivo mediante prompt engineering
  \item Implementé pipeline automatizado de ETL que procesa y limpia datos desde múltiples fuentes, aplicando Machine Learning con Python (scikit-learn) para generar predicciones de tendencias con 87\% de precisión
  \item Diseñé dashboards interactivos con visualizaciones dinámicas en Chart.js y UI/UX optimizada, permitiendo exploración intuitiva de datos y exportación de insights en formatos JSON y CSV
  \item Aseguré integridad de datos mediante validaciones en backend, encriptación de información sensible con bcrypt, y autenticación JWT, cumpliendo estándares de seguridad para aplicaciones productivas
\end{itemize}

\section{Habilidades T\'{e}cnicas}
\textbf{Análisis de Datos \& BI:} SQL Server (Experto en T-SQL, Stored Procedures, Optimización), Power BI (DAX, M, Dashboards), Python (Pandas, NumPy, Matplotlib), Excel Avanzado (Tablas Dinámicas, Macros VBA) \\
\textbf{Bases de Datos \& Cloud:} Azure SQL Database, Oracle, PostgreSQL, MongoDB, Microsoft Azure (Data Factory, Synapse), Modelado Dimensional, ETL \\
\textbf{Desarrollo \& Automatización:} Python (Scripting, Machine Learning con scikit-learn), C\#, Node.js, Git/GitHub, Docker, APIs REST \\
\textbf{IA \& Tecnologías Emergentes:} Prompt Engineering (Claude AI, Gemini), Análisis Predictivo, Data Science Fundamentals \\
\textbf{Frontend \& Visualización:} React, TypeScript, JavaScript (ES6+), HTML5/CSS3, JSON/XML, UI/UX Design

\section{Formaci\'{o}n Acad\'{e}mica}
\textbf{Business Analytics Essentials, Machine Learning con Python, Data Science, Power BI y Optimización de Procesos}, Escuela de Postgrado, Universidad de Chile \hfill 2023-2024 \\
\textbf{Administración de Bases de Datos SQL Server}, Escuela de Postgrado, Universidad de Chile \hfill 2023 \\
\textbf{Desarrollo de Apps Móviles}, Universidad Complutense de Madrid \hfill 2017 \\
\textbf{Analista Programador}, Universidad Tecnológica de Chile INACAP \hfill 2013

\end{document}